\documentclass[10pt]{extarticle}
\usepackage[a4paper,margin=1.5cm]{geometry}
\usepackage[T1]{fontenc}
\usepackage{amsmath,amssymb}
\usepackage{graphicx}
\usepackage[export]{adjustbox}
\usepackage{multicol}
\usepackage{enumitem}
\usepackage{booktabs,array,tabularx}
\setlength{\columnsep}{0.7cm}
\setlength{\parindent}{0pt}
\setlength{\parskip}{3pt}
\pagestyle{empty}
\makeatletter
% figure registry (figdefs.tex is generated by the extract step)
\newcommand{\qfigdef}[4]{\@namedef{qf@file@#1}{#2}\@namedef{qf@w@#1}{#3}\@namedef{qf@wide@#1}{#4}}
\newcommand{\qf@check}[1]{\@ifundefined{qf@file@#1}{\PackageError{score2}{Unknown figure id #1}{}}{}}
\newcommand{\qf@img}[4][]{\includegraphics[width=#2\dimexpr\@nameuse{qf@w@#4}pt\relax,max width=#3,max height=0.45\textheight,#1]{\@nameuse{qf@file@#4}}}
% Full-width questions: finalize wraps every question that contains a wide figure in \qfullbegin...\qfullend,
% so the WHOLE question (stem, figures, options) is set across the page in reading order, then the two columns
% resume. The question's head - its number, stem and first figure (\qfullsplit marks the end of the head) -
% is kept on one page, so a number is never left at a page foot with its figure overleaf; the rest may break
% across pages normally, so no big gaps are left. Without a \qfullsplit, the whole question is kept together
% if it is under 45% of a page. Never write these yourself.
\newif\ifqf@full
\newif\ifqf@boxed
\newsavebox{\qf@box}
\newcommand{\qf@open}{\global\qf@fulltrue\par\global\qf@boxedtrue
  \setbox\qf@box\vbox\bgroup\hsize\textwidth\linewidth\textwidth}
\newcommand{\qf@emit}[1]{\par\egroup\global\qf@boxedfalse
  \ifdim\dimexpr\ht\qf@box+\dp\qf@box\relax<#1\textheight\noindent\usebox{\qf@box}\par\else\unvbox\qf@box\fi}
\newcommand{\qfullsplit}{\ifqf@boxed\qf@emit{0.85}\fi}
\newcommand{\qf@close}{\ifqf@boxed\qf@emit{0.45}\fi\global\qf@fullfalse}
\newcommand{\qfullbegin}{\end{multicols}\qf@open}
\newcommand{\qfullend}{\qf@close\begin{multicols}{2}}
\newcommand{\qfullstart}{\qf@open}
\newcommand{\qfulldone}{\qf@close}
% \qfig[scale]{id}: block figure, centred, directly where it is written (size set by the tool).
% Fallback if a wide figure is not inside a full-width question: break out of the columns in place - no floats.
\newcommand{\qfig}[2][1]{\qf@check{#2}%
  \ifqf@full\qf@blk{#1}{#2}\else
  \ifnum\@nameuse{qf@wide@#2}=1 \expandafter\@firstoftwo\else\expandafter\@secondoftwo\fi
  {\end{multicols}\begin{center}\qf@img{#1}{\textwidth}{#2}\end{center}\begin{multicols}{2}}%
  {\qf@blk{#1}{#2}}\fi}
\newcommand{\qf@blk}[2]{\par\smallskip{\centering\qf@img{#1}{\linewidth}{#2}\par}\smallskip}
% \qfigi[scale]{id}: inline figure (inside option lists / tables / minipages); never breaks columns.
\newcommand{\qfigi}[2][1]{\qf@check{#2}\qf@img[valign=c]{#1}{\linewidth}{#2}}
\makeatother
% structure macros - lint reads these to check completeness and figure placement
\newcommand{\papersec}[1]{\par\bigskip{\centering\Large\bfseries PAPER-#1\par}\medskip}
\newcommand{\partsec}[2]{\par\medskip\noindent{\large\bfseries PART-#1: #2}\par\smallskip}
\newcommand{\sect}[1]{\par\smallskip\noindent\textbf{#1}\par}
\newcommand{\qn}[1]{\par\smallskip\noindent\textbf{#1.}~}
\newcommand{\paraq}[2]{\par\smallskip\noindent\textbf{Paragraph for Questions #1 and #2}\par}
\newenvironment{qoptions}{\begin{itemize}[leftmargin=1.8em,itemsep=1pt,topsep=2pt,label={}]}{\end{itemize}}
\newcommand{\opt}[2]{\item[(#1)] #2}
% figure registry (generated by score2_tools.py finalize)
\qfigdef{p1_phy_s1_q01_1}{Score Advanced - Test 37 (19 May 2024)-img/p1_phy_s1_q01_1.png}{83.5}{0}
\qfigdef{p1_phy_s1_q01_3}{Score Advanced - Test 37 (19 May 2024)-img/p1_phy_s1_q01_3.png}{125.7}{0}
\qfigdef{p1_phy_s1_q04_1}{Score Advanced - Test 37 (19 May 2024)-img/p1_phy_s1_q04_1.png}{117.2}{0}
\qfigdef{p1_phy_s2_q07_1}{Score Advanced - Test 37 (19 May 2024)-img/p1_phy_s2_q07_1.png}{101.8}{0}
\qfigdef{p1_phy_s3_q08_1}{Score Advanced - Test 37 (19 May 2024)-img/p1_phy_s3_q08_1.png}{76.0}{0}
\qfigdef{p1_phy_s3_q10_1}{Score Advanced - Test 37 (19 May 2024)-img/p1_phy_s3_q10_1.png}{86.2}{0}
\qfigdef{p1_phy_s3_q11_P}{Score Advanced - Test 37 (19 May 2024)-img/p1_phy_s3_q11_P.png}{70.3}{0}
\qfigdef{p1_phy_s3_q11_Q}{Score Advanced - Test 37 (19 May 2024)-img/p1_phy_s3_q11_Q.png}{91.8}{0}
\qfigdef{p1_phy_s3_q11_R}{Score Advanced - Test 37 (19 May 2024)-img/p1_phy_s3_q11_R.png}{66.7}{0}
\qfigdef{p1_phy_s3_q11_S}{Score Advanced - Test 37 (19 May 2024)-img/p1_phy_s3_q11_S.png}{96.8}{0}
\qfigdef{p1_phy_s4_q04_1}{Score Advanced - Test 37 (19 May 2024)-img/p1_phy_s4_q04_1.png}{86.2}{0}
\qfigdef{p1_phy_s4_q06_1}{Score Advanced - Test 37 (19 May 2024)-img/p1_phy_s4_q06_1.png}{69.3}{0}
\qfigdef{p1_chem_s1_q03_1}{Score Advanced - Test 37 (19 May 2024)-img/p1_chem_s1_q03_1.png}{165.7}{0}
\qfigdef{p1_chem_s1_q04_1}{Score Advanced - Test 37 (19 May 2024)-img/p1_chem_s1_q04_1.png}{262.0}{1}
\qfigdef{p1_chem_s1_q04_O}{Score Advanced - Test 37 (19 May 2024)-img/p1_chem_s1_q04_O.png}{308.8}{1}
\qfigdef{p1_chem_s2_q05_1}{Score Advanced - Test 37 (19 May 2024)-img/p1_chem_s2_q05_1.png}{102.2}{0}
\qfigdef{p1_chem_s2_q06_1}{Score Advanced - Test 37 (19 May 2024)-img/p1_chem_s2_q06_1.png}{73.0}{0}
\qfigdef{p1_chem_s3_q08_1}{Score Advanced - Test 37 (19 May 2024)-img/p1_chem_s3_q08_1.png}{131.7}{0}
\qfigdef{p1_chem_s3_q10_P}{Score Advanced - Test 37 (19 May 2024)-img/p1_chem_s3_q10_P.png}{107.0}{0}
\qfigdef{p1_chem_s3_q10_Q}{Score Advanced - Test 37 (19 May 2024)-img/p1_chem_s3_q10_Q.png}{93.2}{0}
\qfigdef{p1_chem_s3_q10_R}{Score Advanced - Test 37 (19 May 2024)-img/p1_chem_s3_q10_R.png}{76.3}{0}
\qfigdef{p1_chem_s3_q10_S}{Score Advanced - Test 37 (19 May 2024)-img/p1_chem_s3_q10_S.png}{79.5}{0}
\qfigdef{p1_chem_s4_q05_1}{Score Advanced - Test 37 (19 May 2024)-img/p1_chem_s4_q05_1.png}{241.3}{0}
\qfigdef{p2_phy_s1_q01_1}{Score Advanced - Test 37 (19 May 2024)-img/p2_phy_s1_q01_1.png}{112.2}{0}
\qfigdef{p2_phy_s1_q04_1}{Score Advanced - Test 37 (19 May 2024)-img/p2_phy_s1_q04_1.png}{95.7}{0}
\qfigdef{p2_phy_s2_q05_1}{Score Advanced - Test 37 (19 May 2024)-img/p2_phy_s2_q05_1.png}{85.8}{0}
\qfigdef{p2_phy_s3_para1_2_1}{Score Advanced - Test 37 (19 May 2024)-img/p2_phy_s3_para1_2_1.png}{159.0}{0}
\qfigdef{p2_phy_s3_para3_4_1}{Score Advanced - Test 37 (19 May 2024)-img/p2_phy_s3_para3_4_1.png}{122.6}{0}
\qfigdef{p2_phy_s4_q06_1}{Score Advanced - Test 37 (19 May 2024)-img/p2_phy_s4_q06_1.png}{91.5}{0}
\qfigdef{p2_phy_s4_q07_1}{Score Advanced - Test 37 (19 May 2024)-img/p2_phy_s4_q07_1.png}{76.0}{0}
\qfigdef{p2_phy_s4_q10_1}{Score Advanced - Test 37 (19 May 2024)-img/p2_phy_s4_q10_1.png}{91.5}{0}
\qfigdef{p2_chem_s1_q01_1}{Score Advanced - Test 37 (19 May 2024)-img/p2_chem_s1_q01_1.png}{65.7}{0}
\qfigdef{p2_chem_s1_q01_2}{Score Advanced - Test 37 (19 May 2024)-img/p2_chem_s1_q01_2.png}{42.0}{0}
\qfigdef{p2_chem_s1_q01_3}{Score Advanced - Test 37 (19 May 2024)-img/p2_chem_s1_q01_3.png}{42.2}{0}
\qfigdef{p2_chem_s1_q02_0}{Score Advanced - Test 37 (19 May 2024)-img/p2_chem_s1_q02_0.png}{261.2}{1}
\qfigdef{p2_chem_s1_q02_AL}{Score Advanced - Test 37 (19 May 2024)-img/p2_chem_s1_q02_AL.png}{45.1}{0}
\qfigdef{p2_chem_s1_q02_AP}{Score Advanced - Test 37 (19 May 2024)-img/p2_chem_s1_q02_AP.png}{62.4}{0}
\qfigdef{p2_chem_s1_q02_BL}{Score Advanced - Test 37 (19 May 2024)-img/p2_chem_s1_q02_BL.png}{46.7}{0}
\qfigdef{p2_chem_s1_q02_BP}{Score Advanced - Test 37 (19 May 2024)-img/p2_chem_s1_q02_BP.png}{51.5}{0}
\qfigdef{p2_chem_s1_q02_CL}{Score Advanced - Test 37 (19 May 2024)-img/p2_chem_s1_q02_CL.png}{49.9}{0}
\qfigdef{p2_chem_s1_q02_CP}{Score Advanced - Test 37 (19 May 2024)-img/p2_chem_s1_q02_CP.png}{62.0}{0}
\qfigdef{p2_chem_s1_q02_DL}{Score Advanced - Test 37 (19 May 2024)-img/p2_chem_s1_q02_DL.png}{49.3}{0}
\qfigdef{p2_chem_s1_q02_DP}{Score Advanced - Test 37 (19 May 2024)-img/p2_chem_s1_q02_DP.png}{58.5}{0}
\qfigdef{p2_chem_s1_q03_1}{Score Advanced - Test 37 (19 May 2024)-img/p2_chem_s1_q03_1.png}{138.0}{0}
\qfigdef{p2_chem_s1_q03_A}{Score Advanced - Test 37 (19 May 2024)-img/p2_chem_s1_q03_A.png}{107.8}{0}
\qfigdef{p2_chem_s1_q03_B}{Score Advanced - Test 37 (19 May 2024)-img/p2_chem_s1_q03_B.png}{87.6}{0}
\qfigdef{p2_chem_s1_q03_C}{Score Advanced - Test 37 (19 May 2024)-img/p2_chem_s1_q03_C.png}{75.3}{0}
\qfigdef{p2_chem_s1_q03_D}{Score Advanced - Test 37 (19 May 2024)-img/p2_chem_s1_q03_D.png}{75.0}{0}
\qfigdef{p2_chem_s1_q04_A}{Score Advanced - Test 37 (19 May 2024)-img/p2_chem_s1_q04_A.png}{138.7}{0}
\qfigdef{p2_chem_s1_q04_B}{Score Advanced - Test 37 (19 May 2024)-img/p2_chem_s1_q04_B.png}{141.1}{0}
\qfigdef{p2_chem_s1_q04_C}{Score Advanced - Test 37 (19 May 2024)-img/p2_chem_s1_q04_C.png}{140.4}{0}
\qfigdef{p2_chem_s1_q04_D}{Score Advanced - Test 37 (19 May 2024)-img/p2_chem_s1_q04_D.png}{140.2}{0}
\qfigdef{p2_chem_s3_para1_2_1}{Score Advanced - Test 37 (19 May 2024)-img/p2_chem_s3_para1_2_1.png}{201.3}{0}
\qfigdef{p2_chem_s3_para3_4_1}{Score Advanced - Test 37 (19 May 2024)-img/p2_chem_s3_para3_4_1.png}{243.2}{0}
\qfigdef{p2_chem_s4_q09_1}{Score Advanced - Test 37 (19 May 2024)-img/p2_chem_s4_q09_1.png}{57.6}{0}
\raggedcolumns
\makeatletter\@beginparpenalty=10000\makeatother
\newcolumntype{L}[1]{>{\raggedright\arraybackslash}p{#1}}
\begin{document}
\begin{center}{\Large\bfseries Score Advanced --- Test 37 (19 May 2024)}\end{center}
\begin{multicols}{2}
\raggedcolumns
\papersec{1}
\partsec{1}{PHYSICS}
\sect{SECTION-I(i)}
\qn{1} A block is resting on a horizontal plate in the XY plane and co-efficient of friction between block and the plate is $\mu$. The plate begins to move with velocity v = bt$^2$ in X-direction. At what time will the block starts sliding from plate ?\par
\qfig{p1_phy_s1_q01_1}
\begin{qoptions}
\opt{A}{$\dfrac{\mu b}{g}$}
\opt{B}{$\dfrac{\mu bg}{2}$}
\opt{C}{$\dfrac{\mu g}{b}$}
\opt{D}{$\dfrac{\mu g}{2b}$}
\end{qoptions}
\qn{2} Monochromatic light of wavelength $\lambda$ emerging from slit S illuminates slits S$_1$ and S$_2$ which are placed with respect to S as shown in figure. The distances x and D are large compared to the separation d between the slits. If x = D/2, the minimum value of d so that there is a dark fringe at the centre P of the screen is :-\par
\qfig{p1_phy_s1_q01_3}
\begin{qoptions}
\opt{A}{$\sqrt{\dfrac{\lambda D}{3}}$}
\opt{B}{$\sqrt{\dfrac{2\lambda D}{3}}$}
\opt{C}{$\sqrt{\lambda D}$}
\opt{D}{$2\sqrt{\dfrac{\lambda D}{3}}$}
\end{qoptions}
\qn{3} An electromagnetic wave of frequency f = 3.0 MHz passes from vacuum into a dielectric medium with permittivity e = 4.0. Then :-
\begin{qoptions}
\opt{A}{wavelength is doubled and the frequency remains unchanged}
\opt{B}{wavelength is doubled and frequency becomes half.}
\opt{C}{wavelength is halved and frequency remains unchanged.}
\opt{D}{wavelength and frequency both remain unchanged.}
\end{qoptions}
\qn{4} Switch is in position A for a long time. It is suddenly switched to the side B. Mass of rod is m and length is $\ell$. Maximum speed of rod will be :-\par
\qfig{p1_phy_s1_q04_1}
\begin{qoptions}
\opt{A}{$\left(\dfrac ER\right)\sqrt{\dfrac Lm}$}
\opt{B}{$\left(\dfrac ER\right)\sqrt{\dfrac mL}$}
\opt{C}{$2\left(\dfrac ER\right)\sqrt{\dfrac Lm}$}
\opt{D}{$2\left(\dfrac ER\right)\sqrt{\dfrac mL}$}
\end{qoptions}
\sect{SECTION-I(ii)}
\qn{5} The current density in a wire of radius a varies with radial distance r as J = kr$^2$, where k is a constant.
\begin{qoptions}
\opt{A}{Total current passing through the cross section of the wire is I = $\dfrac{\pi ka^4}2$.}
\opt{B}{Total current passing through the cross section of the wire is I = $\dfrac{3\pi ka^3}2$}
\opt{C}{The field at a distance r $>$ a is B = $\dfrac{\mu_0ka^4}{4r}$}
\opt{D}{The field at a distance r $<$ a is B = $\dfrac{\mu_0kr^3}4$}
\end{qoptions}
\qn{6} When the number of electrons striking the target material in Coolidge's tube (i.e., X-ray tube) is increased keeping the potential difference same, then :-
\begin{qoptions}
\opt{A}{intensity of X-rays increase}
\opt{B}{wavelength of X-rays photons increase}
\opt{C}{frequency of cut off X-rays photons remains same.}
\opt{D}{energy of most energetic X-rays photons increase.}
\end{qoptions}
\qn{7} A smooth track in the form of a quarter circle of radius 6m lies in the vertical plane. A particle moves from P$_1$ to P$_2$ under the action of forces $\vec F_1,\vec F_2$ and $\vec F_3$. Force $\vec F_1$ is always toward P$_2$ and always 20 N in magnitude. Force $\vec F_2$ is always acts horizontally and is always 30 N in magnitude. Force $\vec F_3$ always acts tangentially to the track and is of magnitude 15 N. Select the correct alternative (s) :\par
\qfig{p1_phy_s2_q07_1}
\begin{qoptions}
\opt{A}{Work done by $\vec F_1$ is 120 J}
\opt{B}{Work done by $\vec F_2$ is 180 J}
\opt{C}{Work done by $\vec F_3$ is 45 $\pi$ J}
\opt{D}{$\vec F_1$ is conservative in nature}
\end{qoptions}
\sect{SECTION-I(iii)}
\qn{8} A solid sphere is rotating about an axis as shown in Figure. An insect follows the dotted path on the circumference of sphere as shown. Match the following :-\par
\qfig{p1_phy_s3_q08_1}
{\footnotesize\setlength{\tabcolsep}{1.5pt}\noindent\begin{tabular}{|c|L{3.2cm}|c|L{4cm}|}
\hline
 & \multicolumn{1}{>{\centering\arraybackslash}p{3.2cm}|}{List$-$I} & & \multicolumn{1}{>{\centering\arraybackslash}p{4cm}|}{List$-$II} \\ \hline
(P) & Moment of inertia & (1) & will remain constant \\ \hline
(Q) & Angular velocity & (2) & will first increase then decrease \\ \hline
(R) & Angular momentum & (3) & will first decrease then increases \\ \hline
(S) & Rotational kinetic energy & (4) & will continuously decrease \\ \hline
 & & (5) & will continuously increase \\ \hline
\end{tabular}}\par
\begin{qoptions}
\opt{A}{P $\to$ 2;Q $\to$ 3;R $\to$ 1;S $\to$ 3}
\opt{B}{P $\to$ 1;Q $\to$ 2;R $\to$ 3;S $\to$ 4}
\opt{C}{P $\to$ 5;Q $\to$ 2;R $\to$ 3;S $\to$ 1}
\opt{D}{P $\to$ 3;Q $\to$ 2;R $\to$ 1;S $\to$ 4}
\end{qoptions}
\qn{9} Match the following\par
{\footnotesize\setlength{\tabcolsep}{1.5pt}\noindent\begin{tabular}{|c|L{1.5cm}|c|L{4.5cm}|}
\hline
 & \multicolumn{1}{>{\centering\arraybackslash}p{1.5cm}|}{List$-$I} & & \multicolumn{1}{>{\centering\arraybackslash}p{4.5cm}|}{List$-$II} \\ \hline
(P) & $\dfrac RL$ & (1) & Time \\ \hline
(Q) & CR & (2) & Frequency \\ \hline
(R) & $\dfrac EB$ & (3) & Speed \\ \hline
(S) & $\dfrac1{\mu_0\varepsilon_0}$ & (4) & Kinetic energy \\ \hline
 & & (5) & Kinetic energy per unit mass \\ \hline
\end{tabular}}\par
\begin{qoptions}
\opt{A}{P $\to$ 3;Q $\to$ 2;R $\to$ 5;S $\to$ 1}
\opt{B}{P $\to$ 2;Q $\to$ 1;R $\to$ 3;S $\to$ 5}
\opt{C}{P $\to$ 1;Q $\to$ 2;R $\to$ 3;S $\to$ 4}
\opt{D}{P $\to$ 5;Q $\to$ 3;R $\to$ 1;S $\to$ 4}
\end{qoptions}
\qn{10} With reference to the following P-V diagram, match the following quantities.\par
[W : work done by gas ; $\Delta$U : change in internal energy of gas]\par
\qfig{p1_phy_s3_q10_1}
{\footnotesize\setlength{\tabcolsep}{1.5pt}\noindent\begin{tabular}{|c|L{2.8cm}|c|L{2cm}|}
\hline
 & \multicolumn{1}{>{\centering\arraybackslash}p{2.8cm}|}{List$-$I} & & \multicolumn{1}{>{\centering\arraybackslash}p{2cm}|}{List$-$II} \\ \hline
(P) & Process MN & (1) & W = 0 \\ \hline
(Q) & Process NO & (2) & W $<$ 0 \\ \hline
(R) & Process OP & (3) & W $>$ 0 \\ \hline
(S) & Process PM & (4) & $\Delta$u $>$ 0 \\ \hline
 & & (5) & $\Delta$u $<$ 0 \\ \hline
\end{tabular}}\par
\begin{qoptions}
\opt{A}{P $\to$ 1,3;Q $\to$ 2,4;R $\to$ 1,5;S $\to$ 1,4}
\opt{B}{P $\to$ 1,4;Q $\to$ 3,4;R $\to$ 2,3;S $\to$ 3,4}
\opt{C}{P $\to$ 1,4;Q $\to$ 3,4;R $\to$ 1,5;S $\to$ 2,5}
\opt{D}{P $\to$ 1,2;Q $\to$ 3,4;R $\to$ 1,5;S $\to$ 2,5}
\end{qoptions}
\qn{11} List$-$I shows different charge distributions and short electric dipole at a distance x from the charge distributions. List$-$II gives the dependence of force acting on the dipole as of function of x.\par
{\footnotesize\setlength{\tabcolsep}{1.5pt}\noindent\begin{tabular}{|c|L{4.6cm}|c|L{1.4cm}|}
\hline
 & \multicolumn{1}{>{\centering\arraybackslash}p{4.6cm}|}{List$-$I} & & \multicolumn{1}{>{\centering\arraybackslash}p{1.4cm}|}{List$-$II} \\ \hline
(P) & \qfigi{p1_phy_s3_q11_P}\par Uniformly charged thin disc. & (1) & x$^{-2}$ \\ \hline
(Q) & \qfigi{p1_phy_s3_q11_Q}\par Uniformly charged thick infinite sheet. Dipole is placed inside the sheet at distance x from axis of sheet. & (2) & x$^{-5}$ \\ \hline
(R) & \qfigi{p1_phy_s3_q11_R}\par Uniform infinite line of charge. & (3) & x$^0$ \\ \hline
(S) & \qfigi{p1_phy_s3_q11_S}\par Uniformly charged sphere. Dipole is placed outside the sphere. & (4) & None of these \\ \hline
 & & (5) & x$^{-3}$ \\ \hline
\end{tabular}}\par
\begin{qoptions}
\opt{A}{P $\to$ 4;Q $\to$ 1;R $\to$ 3;S $\to$ 2}
\opt{B}{P $\to$ 4;Q $\to$ 3;R $\to$ 1;S $\to$ 5}
\opt{C}{P $\to$ 1;Q $\to$ 2;R $\to$ 4;S $\to$ 3}
\opt{D}{P $\to$ 4;Q $\to$ 3;R $\to$ 1;S $\to$ 2}
\end{qoptions}
\sect{SECTION-II}
\qn{1} An ammeter and a voltmeter are connected in series to a battery with an emf E = 6V. When a certain resistance is connected in parallel with the voltmeter, the reading of the latter decreases two times, where as the reading of the ammeter increase the same number of times. Find the voltmeter reading, in volt, after the connection of the resistance.
\qn{2} An ideal gas having initial pressure p, volume V and temperature T is allowed to expand adiabatically until its volume becomes 5.66 V while its temperature falls to $\dfrac T2$. Calculate the degrees of freedom of the gas molecules. Given log$_\mathrm e$(5.66) = 1.73, log$_\mathrm e$ (2) = 0.70
\qn{3} The peak emission from a black body at a certain temperature occurs at a wavelength of 6000 \AA. On increasing its temperature, the total radiation emitted is increased 16 times. These radiations are allowed to fall on a metal surface. Photoelectrons emitted by the peak radiation at higher temperature can be bought to rest by applying a potential equivalent to the excitation potential corresponding to the transition for the level n = 4 to n = 2 in the Bohr's hydrogen atom. The work function of the metal is given by $\dfrac\alpha{100}$eV where $\alpha$ is the numerical constant. The value of $\alpha$ is found 53 $\times$ x, find x. [Take : hc = 12420 eV$-$\AA]
\qn{4} The natural frequency of the system shown in figure is f = $\dfrac1\pi\sqrt{\dfrac{xk}M}$. The pulleys are smooth and massless. Find x.\par
\qfig{p1_phy_s4_q04_1}
\qn{5} A road passes at some distance from a standing man. A truck is coming on the road with some acceleration. The truck driver blows a whistle of frequency 500 Hz when the line joining the truck and the man makes an angle of $\theta$ with the road. The man hears a note having a frequency of 600 Hz when the truck is closest to him. Also, the speed of truck has got doubled during this time. Find the value of $\theta$, in degree.
\qn{6} A uniform disc of radius r and mass m is charged uniformly with the charge q. This disc is placed flat on a rough horizontal surface having coefficient of friction $\mu$. A uniform magnetic field is present in a circular region (a $>$ r) but varying as kt$^3$ as shown in Figure. Find the time, in second after which the disc begins to rotate. (Given r = 1m, m = 18 kg, q = 1 C, $\mu$ = 0.1, K = 4, g = 10 ms$^{-2}$)\par
\qfig{p1_phy_s4_q06_1}
\partsec{2}{CHEMISTRY}
\sect{SECTION-I(i)}
\qn{1} [X] $\mathop{\rightleftharpoons}\limits^{\mathrm{OH}^-}_{\mathrm H^+}$ [Y]\par
Both [X] and [Y] produce ink blue solution with H$_2$O$_2$ in presence of organic solvents and H$_2$SO$_4$, [X] \& [Y] respectively
\begin{qoptions}
\opt{A}{MnO$_4^-$, MnO$_4^{-2}$}
\opt{B}{Cr$_2$O$_7^{-2}$, CrO$_4^{-2}$}
\opt{C}{CrO$_4^{-2}$, Cr$_2$O$_7^{-2}$}
\opt{D}{SO$_3^{-2}$, SO$_4^{-2}$}
\end{qoptions}
\qn{2} Which of the following set contains, all tetrahedral geometries with respect to central metal atom or ion ?
\begin{qoptions}
\opt{A}{FeCl$_4^{2-}$, Cu(CN)$_4^{3-}$, PdCl$_4^{2-}$}
\opt{B}{NiCl$_2$(PPh$_3$)$_2$, Pt(CN)$_4^{2-}$, RhCl(PPh$_3$)$_3$}
\opt{C}{Ni(CO)$_4$, Fe(CO)$_4^{2-}$, Cu(CN)$_4^{3-}$}
\opt{D}{Fe(CO)$_5$ , Co(CO)$_4^-$ , [NiCl$_4$]$^{2-}$}
\end{qoptions}
\qn{3} A small amount of solid lead iodide was added to a beaker of water, which was stirred. Most of the solid settled at the bottom of the beaker, but a little dissolved, establishing the equilibrium\par
PbI$_2$ (s) $\rightleftharpoons$ Pb$^{2+}$ (aq) + 2I$^-$ (aq)\par
The rates of the forward and reverse reactions were monitored over time, producing the graph shown below:\par
\qfig{p1_chem_s1_q03_1}
What happened at time t?
\begin{qoptions}
\opt{A}{The beaker was cooled in an ice-bath.}
\opt{B}{A small amount of solid KI was added to the beaker.}
\opt{C}{A small amount of solid PbI$_2$ was removed from the beaker.}
\opt{D}{A small amount of water was added to the beaker.}
\end{qoptions}
\qfullbegin
\qn{4} Consider the following scheme:\par
\qfig{p1_chem_s1_q04_1}
\qfullsplit
If P is ethyl chloride, R can be\par
\qfig{p1_chem_s1_q04_O}
\begin{qoptions}
\opt{A}{I, II, III, IV}
\opt{B}{II, III, IV only}
\opt{C}{IV only}
\opt{D}{II, IV only}
\end{qoptions}
\qfullend
\sect{SECTION-I(ii)}
\qn{5} One mole of a diatomic ideal gas is subjected to a reversible process that involves two steps AB and AC. The pressure at A and C is same. Consider the graph and choose correct statements.\par
[Given : $\ell$n 2 = 0.70, $\ell$n 3 = 1.1, R = 0.082 L-atm/K-mol]\par
\qfig{p1_chem_s2_q05_1}
\begin{qoptions}
\opt{A}{Entropy change for process AC is 0.2 kJK$^{-1}$}
\opt{B}{In Path AB work done is $-$22.3 J}
\opt{C}{Volume of gas at B = 2 $\times$ Volume of gas at A}
\opt{D}{Volume of gas at C is 32.8 lit.}
\end{qoptions}
\qn{6} The correct statement/s about the pyranose form of a sugar (X) given below is/are:\par
\qfig{p1_chem_s2_q06_1}
\begin{qoptions}
\opt{A}{It exists in two anomeric pyranose forms}
\opt{B}{It reacts with Tollen's reagent to give a silver mirror}
\opt{C}{The penta-O-methyl derivative of (X) is non reducing}
\opt{D}{It resists reduction with aqueous sodium borohydride}
\end{qoptions}
\qn{7} The correct statement(s) related to oxoacids of phosphorous is(are) :
\begin{qoptions}
\opt{A}{Upon heating, H$_3$PO$_3$ undergoes disproportionation reaction to produce H$_3$PO$_4$ and PH$_3$.}
\opt{B}{While H$_3$PO$_3$ can act as reducing agent, H$_3$PO$_4$ cannot.}
\opt{C}{H$_3$PO$_3$ is a monobasic acid.}
\opt{D}{The H atom of P-H bond in H$_3$PO$_3$ is not ionizable in water.}
\end{qoptions}
\sect{SECTION-I(iii)}
\qn{8} Study the sequence of reactions of diborane (B$_2$H$_6$) given below and match the products K, L, M and N with products mention in List-II (shown reagents are taken in excess)\par
\qfig{p1_chem_s3_q08_1}
{\footnotesize\setlength{\tabcolsep}{1.5pt}\noindent\begin{tabular}{|c|L{1cm}|c|L{3.5cm}|}
\hline
 & \multicolumn{1}{>{\centering\arraybackslash}p{1cm}|}{List-I} & & \multicolumn{1}{>{\centering\arraybackslash}p{3.5cm}|}{List-II} \\ \hline
(P) & K & (1) & [BH$_2$(NH$_3$)$_2$][BH$_4$] \\ \hline
(Q) & L & (2) & B$_2$H$_5$Cl \\ \hline
(R) & M & (3) & BCl$_3$ \\ \hline
(S) & N & (4) & B$_3$N$_3$H$_3$Cl$_3$ \\ \hline
 & & (5) & B$_3$N$_3$H$_6$ \\ \hline
\end{tabular}}\par
\begin{qoptions}
\opt{A}{P $\to$ 5;Q $\to$ 1;R $\to$ 3;S $\to$ 2}
\opt{B}{P $\to$ 5;Q $\to$ 1;R $\to$ 2;S $\to$ 4}
\opt{C}{P $\to$ 1;Q $\to$ 5;R $\to$ 2;S $\to$ 4}
\opt{D}{P $\to$ 1;Q $\to$ 5;R $\to$ 3;S $\to$ 4}
\end{qoptions}
\qn{9} Match the half-life of the reaction in List-I with the integrated rate equation in List-II\par
(a = initial concentration of reactant and a $-$ x $-$ concentration of reactant at time `t'.)\par
{\footnotesize\setlength{\tabcolsep}{1.5pt}\noindent\begin{tabular}{|c|L{2.3cm}|c|L{3.6cm}|}
\hline
 & \multicolumn{1}{>{\centering\arraybackslash}p{2.3cm}|}{List $-$ I} & & \multicolumn{1}{>{\centering\arraybackslash}p{3.6cm}|}{List $-$ II} \\ \hline
(P) & $t_{1/2}=\dfrac{0.586}{k}\sqrt a$ & (1) & $k=\dfrac xt$ \\ \hline
(Q) & $t_{1/2}=\dfrac{0.693}{k}$ & (2) & $k=\dfrac1t\left[\dfrac1{a-x}-\dfrac1a\right]$ \\ \hline
(R) & $t_{1/2}=\dfrac a{2k}$ & (3) & $k=\dfrac2t\left[\sqrt a-\sqrt{a-x}\right]$ \\ \hline
(S) & $t_{1/2}=\dfrac1{ka}$ & (4) & $k=\dfrac{2.303}t\log\dfrac a{a-x}$ \\ \hline
 & & (5) & $k=\dfrac2t\left[\dfrac1{\sqrt{a-x}}-\dfrac1{\sqrt a}\right]$ \\ \hline
\end{tabular}}\par
\begin{qoptions}
\opt{A}{P $\to$ 5;Q $\to$ 4;R $\to$ 1;S $\to$ 2}
\opt{B}{P $\to$ 3;Q $\to$ 4;R $\to$ 1;S $\to$ 2}
\opt{C}{P $\to$ 2;Q $\to$ 3;R $\to$ 4;S $\to$ 2}
\opt{D}{P $\to$ 3;Q $\to$ 2;R $\to$ 4;S $\to$ 5}
\end{qoptions}
\qn{10} Match each species in List - I with appropriate statement (s) of List - II and choose correct option\par
{\footnotesize\setlength{\tabcolsep}{1.5pt}\noindent\begin{tabular}{|c|L{3.7cm}|c|L{3cm}|}
\hline
 & \multicolumn{1}{>{\centering\arraybackslash}p{3.7cm}|}{List - I} & & \multicolumn{1}{>{\centering\arraybackslash}p{3cm}|}{List - II} \\ \hline
(P) & \qfigi{p1_chem_s3_q10_P} & (1) & At least one of the product shows optical isomerism. \\ \hline
(Q) & \qfigi{p1_chem_s3_q10_Q} & (2) & Major product contains even number of $\alpha$ - hydrogen's (Hydrogen's which can involve in hyper conjugation). \\ \hline
(R) & \qfigi{p1_chem_s3_q10_R} & (3) & One of the product shows geometrical isomerism. \\ \hline
(S) & \qfigi{p1_chem_s3_q10_S} & (4) & Reactant can turn acidified dichromate to green. \\ \hline
 & & (5) & Will not undergo dehydration. \\ \hline
\end{tabular}}\par
\begin{qoptions}
\opt{A}{P $\to$ 2,3;Q $\to$ 2,4,5;R $\to$ 1,4;S $\to$ 5}
\opt{B}{P $\to$ 2,3,4;Q $\to$ 2,3,4;R $\to$ 1,2;S $\to$ 3,4}
\opt{C}{P $\to$ 1,2,3,4;Q $\to$ 2,3,4;R $\to$ 1,3;S $\to$ 2}
\opt{D}{P $\to$ 2,3,4;Q $\to$ 1,2,3,4;R $\to$ 1,4;S $\to$ 5}
\end{qoptions}
\qn{11} {\footnotesize\setlength{\tabcolsep}{1.5pt}\noindent\begin{tabular}{|c|L{3.6cm}|c|L{2.8cm}|}
\hline
 & \multicolumn{1}{>{\centering\arraybackslash}p{3.6cm}|}{List - I} & & \multicolumn{1}{>{\centering\arraybackslash}p{2.8cm}|}{List - II} \\ \hline
(P) & A compound that liberates ammonia on treatment with hot aqueous KOH & (1) & CH$_3$$-$CH$_2$$-$CN \\ \hline
(Q) & A compound that gives CO$_2$, NH$_4^+$ and acetic acid after boiling in concentrated, aqueous HCl & (2) & CH$_3$$-$CH$_2$$-$CONH$_2$ \\ \hline
(R) & A compound that gives NH$_4^+$ and mono carboxylic acid after treatment with concentrated aqueous HCl. & (3) & NC$-$CH$_2$$-$CN \\ \hline
(S) & A compound that gives product A with MeMgI followed by hydrolysis, which on further treatment with MeMgI excess followed by acidification gives 2-methyl-2-butanol & (4) & CH$_3$COONH$_4$ \\ \hline
 & & (5) & PhCOONH$_4$ \\ \hline
\end{tabular}}\par
The correct option is\_\_\_\_\_\_.
\begin{qoptions}
\opt{A}{P $\to$ 1;Q $\to$ 3;R $\to$ 5;S $\to$ 1}
\opt{B}{P $\to$ 1;Q $\to$ 2;R $\to$ 3;S $\to$ 4}
\opt{C}{P $\to$ 3;Q $\to$ 4;R $\to$ 5;S $\to$ 4}
\opt{D}{P $\to$ 2;Q $\to$ 4;R $\to$ 2;S $\to$ 1}
\end{qoptions}
\sect{SECTION-II}
\qn{1} Find the number of reactions, which give N$_2$.\par
(I) (NH$_4$)$_2$Cr$_2$O$_7\xrightarrow{\Delta}$ \quad (II) Ba(N$_3$)$_2\xrightarrow{\Delta}$\par
(III) NH$_4$Cl $\xrightarrow{\Delta}$ \quad (IV) NH$_3$ (excess) + Cl$_2$ $\rightarrow$\par
(V) NH$_4$NO$_3$ $\xrightarrow[250^\circ\mathrm C]{\Delta}$ \quad (VI) (NH$_4$)$_2$SO$_4$ $\xrightarrow[\mathrm{aq.}]{\mathrm{KOH}}$\par
(VII) C$_2$H$_5$NH$_2$ + HNO$_2$ $\xrightarrow{\mathrm{water}}$
\qn{2} If `x' is the total number of p$\pi$ $-$d$\pi$ bonds and `y' is the total number of lone pair(s) on non terminal atoms per molecule in trimer of SO$_3$, the value of (x + y) is :-
\qn{3} One gram of commercial AgNO$_3$ is dissolved in 50 mL of water. It is treated with 50 mL of a KI solution. The silver iodide thus precipitated is filtered off. Excess of KI in the filtrate is titrated with M/10 KIO$_3$ solution in presence of 6 M HCl till all I$^-$ ions are converted into ICl. It requires 50 mL of M/10 KIO$_3$ solution. Twenty milliliters of the same stock solution of KI requires 30 mL of M/10 KIO$_3$ under similar conditions. The percentage of AgNO$_3$ in the sample is (Atomic weight of silver = 108)
\qn{4} Both gold and silver are extracted by cyanide process. Gold in nature is frequently alloyed with silver, which is also oxidized by aerated NaCN solution, similar to gold. 1000 L of a solution of 0.002 M [Au(CN)$_2$]$^-$ and 0.003 M [Ag(CN)$_2$]$^-$ was evaporated to half of the original volume and treated with 78 g zinc (atomic mass = 65).\par
Assume that elevation from standard conditions is unimportant and that the redox reactions go to completion. Given:\par
[Zn(CN)$_4$]$^{2-}$ + 2e$^-$ $\to$ Zn (s) + 4CN$^-$; E$^\circ$ = $-$1.26V\par
[Au(CN)$_2$]$^-$ + e$^-$ $\to$ Au (s) + 2CN$^-$; E$^\circ$ = $-$0.60V\par
[Ag(CN)$_2$]$^-$ + e$^-$ $\to$ Ag (s) + 2CN$^-$; E$^\circ$ = $-$0.31V\par
If the concentration of [Au(CN)$_2$]$^-$in the solution after addition of zinc in it, be x mol/lit. The value of (x $\times$ 10$^5$) is.
\qn{5} Consider the following reactions\par
\qfig{p1_chem_s4_q05_1}
Amount (in mg) of B reacts completely with 100 ml of 0.1 M NaHCO$_3$ solution.
\qn{6} The correct order of the given reagents to convert benzene to m-Chlorobromobenzene is\par
(1) Sulfuric acid (conc.) and heat\par
(2) Cl$_2$ + FeCl$_3$ and heat\par
(3) NaNO$_3$ + H$_3$O$^+$, 0$^\circ$C\par
(4) H$_2$/Pt catalyst\par
(5) Mg in ether\par
(6) PBr$_3$\par
(7) H$_3$PO$_2$ (aqueous)\par
(8) HNO$_3$ (conc.) + H$_2$SO$_4$ (conc.) and heat\par
(9) Cu$_2$Br$_2$ +HBr\par
(10) (CH$_3$CO)$_2$ O + pyridine\par
Write the correct order in form of their number and express last two digits as answer.\par
[ For, e.g. if order is 1357, then your answer will be 57]
\partsec{3}{MATHEMATICS}
\sect{SECTION-I(i)}
\qn{1} A ray of light travelling along the line OP (O being origin) is reflected by the line mirror 2x $-$ 3y + 1 = 0, the point of incidence being P(1,1). The reflected ray, travelling along PQ is again reflected by the line mirror 2x $-$ 3y $-$ 1 = 0, the point of incidence being Q. From Q, ray move along QR, where R lies on the line 2x $-$ 3y + 1 = 0. Then sum of ordinate of point R and abscissa of point Q is :
\begin{qoptions}
\opt{A}{$\dfrac{100}{13}$}
\opt{B}{$\dfrac{73}{13}$}
\opt{C}{$\dfrac{53}{13}$}
\opt{D}{$\dfrac{23}{13}$}
\end{qoptions}
\qn{2} Number of integers in the range of `m' for which point (m,1) lies in the smaller region bounded by x$^2$ + y$^2$ $-$ 3x + 1 = 0 and 2x $-$ y = 2 is-
\begin{qoptions}
\opt{A}{0}
\opt{B}{1}
\opt{C}{2}
\opt{D}{3}
\end{qoptions}
\qn{3} Let $g(x)=\displaystyle\int_{x^2}^{x^3}\dfrac{dt}{\ln t}$ for x $>$ 1 and $h(x)=\displaystyle\int_1^x\left(2t^2-\ln t\right)g(t)dt$ (x $>$ 1), then :
\begin{qoptions}
\opt{A}{h is increasing on (1, $\infty$)}
\opt{B}{h is decreasing on (1, 2) and increasing on (2, $\infty$)}
\opt{C}{h is decreasing on (1, $\infty$)}
\opt{D}{h is increasing on (1, 2) and decreasing on (2, $\infty$)}
\end{qoptions}
\qn{4} The tangent at (x$'$, y$'$) to the curve x$^3$ + y$^3$ = a$^3$ meets the curve again at (x$_1$, y$_1$), then $\left|\dfrac{x_1}{x'}+\dfrac{y_1}{y'}\right|^2=$
\begin{qoptions}
\opt{A}{4}
\opt{B}{1}
\opt{C}{a}
\opt{D}{a$^2$}
\end{qoptions}
\sect{SECTION-I(ii)}
\qn{5} Let f(x) = sin x + px + q. Then f(x) = 0 has :
\begin{qoptions}
\opt{A}{Only one real root which is positive if p $>$ 1, q $<$ 0}
\opt{B}{Only one real root which is negative if p $<$ $-$1, q $<$ 0}
\opt{C}{Only one real root which is negative if p $>$ 1, q $>$ 0}
\opt{D}{Only one real root which is positive if p $<$ $-$ 1, q $<$ 0}
\end{qoptions}
\qn{6} Let $\vec u=\cos\alpha\hat i+\sin\alpha\hat j$, $\vec v=\cos\beta\hat i+\sin\beta\hat j$, $\vec w=\cos\gamma\hat i+\sin\gamma\hat j$, $0\leq\alpha,\beta,\gamma<2\pi$. If $\vec u.\vec w=\vec v.\vec w=-\sqrt3\vec u.\vec v$ then $\left(2\vec u.\vec v\right)^3$ is equal to :
\begin{qoptions}
\opt{A}{$\dfrac{-8}{27}$}
\opt{B}{$\dfrac{-1}8$}
\opt{C}{1}
\opt{D}{$\dfrac19$}
\end{qoptions}
\qn{7} P is a non$-$singular matrix of order n $\times$ n such that 3PQP$^{-1}$ + P = 2P$^{-1}$ QP then
\begin{qoptions}
\opt{A}{P \& Q both are identity matrices}
\opt{B}{det.(P + Q) = 0}
\opt{C}{PQP$^{-1}$ $-$ P$^{-1}$QP is a singular matrix}
\opt{D}{P + Q is non - invertible matrix}
\end{qoptions}
\sect{SECTION-I(iii)}
\qn{8} Matching the following statement\par
{\footnotesize\setlength{\tabcolsep}{1.5pt}\noindent\begin{tabular}{|c|L{5.2cm}|c|L{0.9cm}|}
\hline
 & \multicolumn{1}{>{\centering\arraybackslash}p{5.2cm}|}{List-I} & & \multicolumn{1}{>{\centering\arraybackslash}p{0.9cm}|}{List-II} \\ \hline
(P) & If $\dfrac{dy}{dx}\left(\dfrac{1+\cos x}y\right)=-\sin x$ and f$\left(\dfrac\pi2\right)=-1$ then f(0) is & (1) & 5 \\ \hline
(Q) & The number of points of non-derivability of the function f(x) = $\left[\dfrac{2x}\pi\right]$ sgn$\left(\dfrac1{\{x\}}\right)$ in ($-$2, 2) is (where [P], \{P\} \& sgn(P) denotes greater integer $\leq$ P, fractional part of P \& signum function of P respectively) & (2) & $-$2 \\ \hline
(R) & The area bounded by $e^{\ln|x+1|}\geqslant|y|,|x|\leqslant1$ is & (3) & 4 \\ \hline
(S) & Number of real values of x which satisfy the equation $\dfrac{\left(\sec^{-1}x\right)^3+\left(\mathrm{cosec}^{-1}x\right)^3}{\left(\tan^{-1}x+\cot^{-1}x\right)^3}=7$ is: & (4) & 1 \\ \hline
 & & (5) & 7 \\ \hline
\end{tabular}}\par
\begin{qoptions}
\opt{A}{P $\to$ 1;Q $\to$ 2;R $\to$ 3;S $\to$ 4}
\opt{B}{P $\to$ 2;Q $\to$ 1;R $\to$ 3;S $\to$ 4}
\opt{C}{P $\to$ 2;Q $\to$ 3;R $\to$ 1;S $\to$ 4}
\opt{D}{P $\to$ 2;Q $\to$ 1;R $\to$ 4;S $\to$ 3}
\end{qoptions}
\qn{9} Matching the following statement\par
{\footnotesize\setlength{\tabcolsep}{1.5pt}\noindent\begin{tabular}{|c|L{5.2cm}|c|L{0.9cm}|}
\hline
 & \multicolumn{1}{>{\centering\arraybackslash}p{5.2cm}|}{List-I} & & \multicolumn{1}{>{\centering\arraybackslash}p{0.9cm}|}{List-II} \\ \hline
(P) & If sin$^2$ A = x and sinA $\cdot$ sin2A sin3A $\cdot$ sin4A = ax$^2$ + bx$^3$ + cx$^4$ + dx$^5$, then the value of a + b + c + d + 5 is & (1) & 1 \\ \hline
(Q) & The values of x between 0 and 2$\pi$ which satisfy the equation $\sin x\sqrt{8\cos^2x}=1$ are in A.P. with common difference d, then $\dfrac{4d}\pi$ is & (2) & 2 \\ \hline
(R) & The least difference between the roots, in the first quadrant $\left(0\leqslant x\leqslant\dfrac\pi2\right)$, of the equation $4\cos x\left(2-3\sin^2x\right)+(\cos2x+1)=0$ is $\dfrac\pi k$, then k = & (3) & 6 \\ \hline
(S) & Let $f(x)=\sqrt{\sin^4x+4\cos^2x}-\sqrt{\cos^4x+4\sin^2x}$ and g(sin 2t) = sint + cos t $\forall$ t $\in\left[\dfrac{-\pi}4,\dfrac\pi4\right]$. If the range of g(f(x)) is [a, b], then the value of a$^2$ + b$^2$ is\_\_\_\_\_\_. & (4) & 5 \\ \hline
 & & (5) & 4 \\ \hline
\end{tabular}}\par
\begin{qoptions}
\opt{A}{P $\to$ 2;Q $\to$ 4;R $\to$ 1;S $\to$ 3}
\opt{B}{P $\to$ 4;Q $\to$ 1;R $\to$ 2;S $\to$ 3}
\opt{C}{P $\to$ 2;Q $\to$ 1;R $\to$ 4;S $\to$ 3}
\opt{D}{P $\to$ 4;Q $\to$ 1;R $\to$ 3;S $\to$ 2}
\end{qoptions}
\qn{10} Matching the following statement\par
{\footnotesize\setlength{\tabcolsep}{1.5pt}\noindent\begin{tabular}{|c|L{5.2cm}|c|L{0.9cm}|}
\hline
 & \multicolumn{1}{>{\centering\arraybackslash}p{5.2cm}|}{List- I} & & \multicolumn{1}{>{\centering\arraybackslash}p{0.9cm}|}{List-II} \\ \hline
(P) & If $\displaystyle\lim_{h\to0^+}\int_{-1}^1\dfrac{h\,dx}{h^2+x^2}=k\pi$ Then k equals & (1) & 5 \\ \hline
(Q) & The area of the region defined by 1 $\leqslant$ $|x-2|+|y+1|$ $\leqslant$ 2 is & (2) & 3 \\ \hline
(R) & $\displaystyle\int_0^{110\pi}[\sin x+\cos x]\,dx=-11k\pi$ then k equal to (where [.] denotes greatest integer function) & (3) & 1 \\ \hline
(S) & If distance between two focii of ellipse $(x-2y+3)^2+(8x+4y+4)^2=20$ is k then $\dfrac{k^2}5$ equals to & (4) & 6 \\ \hline
 & & (5) & 9 \\ \hline
\end{tabular}}\par
\begin{qoptions}
\opt{A}{P $\to$ 3;Q $\to$ 4;R $\to$ 1;S $\to$ 2}
\opt{B}{P $\to$ 2;Q $\to$ 3;R $\to$ 1;S $\to$ 4}
\opt{C}{P $\to$ 3;Q $\to$ 4;R $\to$ 2;S $\to$ 1}
\opt{D}{P $\to$ 2;Q $\to$ 4;R $\to$ 1;S $\to$ 3}
\end{qoptions}
\qn{11} {\footnotesize\setlength{\tabcolsep}{1.5pt}\noindent\begin{tabular}{|c|L{5.2cm}|c|L{0.9cm}|}
\hline
 & \multicolumn{1}{>{\centering\arraybackslash}p{5.2cm}|}{List- I} & & \multicolumn{1}{>{\centering\arraybackslash}p{0.9cm}|}{List-II} \\ \hline
(P) & If $\sqrt{x^2+y^2}+\sqrt{(x+3)^2+(y+4)^2}=9-a^2$ represent an ellipse then number of integral value of a is & (1) & 7 \\ \hline
(Q) & If a, b, c $>$ 0 \& a$^2$ = bc \& a + b + c = abc then the least possible value of a$^2$ is & (2) & 5 \\ \hline
(R) & The minimum value of n $\in$ N so that the sum of coefficient of terms in the expansion of (a +b + c $-$ d)$^n$ is more than the number of terms is:- & (3) & 3 \\ \hline
(S) & If $\displaystyle\lim_{x\to\infty}\left(8x^3+ax^2\right)^{\frac13}-bx$ exists \& is equal to 1, then value of a $-$ 3b is equals to & (4) & 6 \\ \hline
 & & (5) & 2 \\ \hline
\end{tabular}}\par
\begin{qoptions}
\opt{A}{P $\to$ 3;Q $\to$ 4;R $\to$ 1;S $\to$ 1}
\opt{B}{P $\to$ 2;Q $\to$ 1;R $\to$ 4;S $\to$ 1}
\opt{C}{P $\to$ 2;Q $\to$ 1;R $\to$ 4;S $\to$ 4}
\opt{D}{P $\to$ 3;Q $\to$ 3;R $\to$ 1;S $\to$ 4}
\end{qoptions}
\sect{SECTION-II}
\qn{1} The solution of the equation $\left(2x^2+3y^2-7\right)x\,dx-\left(3x^2+2y^2-8\right)y\,dy=0$ is given by $\left(x^2+y^2+\gamma\right)=C\left(x^2-y^2+\delta\right)^\beta$. Where C is arbitrary constant, then ($\beta$ + $\gamma$ + $\delta$) equals
\qn{2} If z is a complex number satisfying the condition $\left|z-\dfrac5z\right|=4$, then sum of greatest and least value of $|z|$ is
\qn{3} If T$_\mathrm n$ is term which is independent of `x' in the expansion of $\left(\dfrac{x^4+x^2+1}{x^2-x+1}-\dfrac{x^2+x+1}x\right)^{12}$, then n is
\qn{4} In a drawer Mr. X has 5 pairs of socks, each pair a different color. On Monday, he selects two individual socks at random from the 10 socks in the drawer. On Tuesday, he selects 2 of the remaining 8 socks at random and on Wednesday, two of the remaining 6 socks at random. If the probability that Wednesday is the first day he selects matching socks is $\dfrac mn$, where m and n are relatively prime positive integers, then the value of $\dfrac{m+n}{11}$ is equal to
\qn{5} A 3 $\times$ 3 determinant has entries either 1 or $-$1. Let S be the set of all determinants such that the product of elements of any row or column is $-$1. For example $\begin{vmatrix}1&-1&1\\1&1&-1\\-1&1&1\end{vmatrix}$ is an element of S and number of elements of S is m. Let P = $\begin{vmatrix}3&-2&3\\2&-2&3\\0&-1&1\end{vmatrix}$ and trace of the matrix adj (adj P) is n then the value of $\dfrac mn$ is
\qn{6} Let the area of given triangle ABC be $\Delta$. Points A$_1$, B$_1$ and C$_1$ are the mid points of the BC, CA and AB respectively. Point A$_2$ is the mid point of CA$_1$. Lines C$_1$A$_1$ and AA$_2$ meet the median BB$_1$ at points E and D respectively. If $\Delta_1$ be the area of the quadrilateral A$_1$A$_2$DE. If the value of $\dfrac{\Delta_1}\Delta=\dfrac mn$ (where m and n are coprime), then n $-$ 3m = \_\_\_\_\_.
\raggedcolumns
\papersec{2}
\partsec{1}{PHYSICS}
\sect{SECTION-I(i)}
\qn{1} A horizontal tube of uniform cross-sectional area A is bent in the form of U as shown in figure. If the liquid of density $\rho$ enters and leaves the tube with velocity v, then the external force F required to hold the bend stationary is :-\par
\qfig{p2_phy_s1_q01_1}
\begin{qoptions}
\opt{A}{F = 0}
\opt{B}{$\rho$Av$^2$}
\opt{C}{2$\rho$Av$^2$}
\opt{D}{None of these}
\end{qoptions}
\qn{2} The activity of a radioactive sample goes down to about 6\% in a time of 2 hours. The half life of sample in minutes is about :-
\begin{qoptions}
\opt{A}{30}
\opt{B}{15}
\opt{C}{60}
\opt{D}{120}
\end{qoptions}
\qn{3} A planet of mass m is revolving in an elliptical orbit about the sun with an orbital period T. If M$_\mathrm{sun}$ be the mass of sun and A be the area of orbit, then the angular momentum of the planet is :-
\begin{qoptions}
\opt{A}{$\dfrac{\mathrm{mA}}{2\mathrm T}$}
\opt{B}{$\dfrac{\mathrm{mA}}{4\mathrm T}$}
\opt{C}{$\dfrac{2\mathrm{mA}}{\mathrm T}$}
\opt{D}{$\dfrac{4\mathrm{mA}}{\mathrm T}$}
\end{qoptions}
\qn{4} Figure shows a thin metal sheet in the plane y = 0, for which the current of constant density flows in the positive x-direction. It is in a constant homogeneous magnetic field of value B$_0\hat{\mathrm k}$. As a result of superposition of magnetic field in region y $>$ 0, the induction field B$_1\hat{\mathrm k}$ and in y $<$ 0 is B$_2\hat{\mathrm k}$ where B$_1 >$ B$_2$. Specify the correct statement :-\par
\qfig{p2_phy_s1_q04_1}
\begin{qoptions}
\opt{A}{B$_0 = \dfrac{(\mathrm B_1-\mathrm B_2)}2$}
\opt{B}{B$_0 = \dfrac{(\mathrm B_1+\mathrm B_2)}2$}
\opt{C}{B$_0$ = B$_1$ + B$_2$}
\opt{D}{B$_0$ = B$_1$ $-$ B$_2$}
\end{qoptions}
\sect{SECTION-I(ii)}
\qn{5} Consider a thermodynamics cycle in a PV diagram shown in the figure performed by one mole of a monoatomic gas. The temperature at A is T$_0$ and volumes are related as V$_\mathrm B$ = V$_\mathrm C$ = 2V$_\mathrm A$. Choose the correct options :-\par
\qfig{p2_phy_s2_q05_1}
\begin{qoptions}
\opt{A}{The maximum temperature during the cycle is 4T$_0$}
\opt{B}{Net work done by the gas during the cycle is 0.5 RT$_0$}
\opt{C}{The heat capacity of the process AB is 2R.}
\opt{D}{The efficiency of the cycle is 8.33\%}
\end{qoptions}
\qn{6} The plates of a parallel plate capacitor with no dielectric are connected to a voltage source. Now a dielectric of dielectric constant K is inserted to fill the whole space between the plates with voltage source remaining connected to the capacitor. Select the correct statement(s) :-
\begin{qoptions}
\opt{A}{The energy stored in the capacitor will become K-times the initial energy.}
\opt{B}{The electric field inside the capacitor will decrease K-times the initial field.}
\opt{C}{The force of attraction between the plates will become K$^2$-times the initial force.}
\opt{D}{The charge on the capacitor will become K-times the initial charge.}
\end{qoptions}
\qn{7} In a resonance tube experiment to determine speed of sound, air column in the pipe is made to resonate with a given fork of frequency 480 Hz. At room temperature T using a resonance tube a student finds two resonance position 16.0 cm and 51.0 cm, length of the air column for the first and the second resonance respectively. It is known that speed of the sound at 273 K is equal to 330 m/s. Then,
\begin{qoptions}
\opt{A}{percentage error in the measurement of speed of sound is $\dfrac47$\%}
\opt{B}{diameter of the tube is 5 cm.}
\opt{C}{temperature of the room T is nearly 283 K.}
\opt{D}{percentage error in the determination of T is $\dfrac87$\%}
\end{qoptions}
\sect{SECTION-II(i)}
\paraq{1}{2}
Two wave pulses are travelling in opposite direction with speed 1 ms$^{-1}$. Figure shows the shape of pulse at t = 0. Based on the above facts, answer the following questions.\par
\qfig{p2_phy_s3_para1_2_1}
\qn{1} The speed of particle (in m/s) at x = 2 cm and t = 0 is :-
\qn{2} The displacement of particle (in mm) at x = 8 cm and t = 0.06 sec is
\paraq{3}{4}
A small block of mass m = 1kg is placed over a wedge of mass M = 4kg as shown in figure. Mass m is released from rest. All surfaces are smooth and the origin O is fixed on the ground shown. Based on above information, answer the following questions.\par
\qfig{p2_phy_s3_para3_4_1}
\qn{3} Final velocity of the wedge is (in m/s)
\qn{4} The block strikes the x-axis at (in m)
\sect{SECTION-II(ii)}
\qn{5} In LCR circuit current resonant frequency is 600 Hz and half power points are at 650 and 550 Hz. The quality factor is
\qn{6} In a zero gravity region, a drop of a liquid of surface tension $\sigma$ assumes a cylindrical shape of diameter D between two parallel glass plates that are a distance `d' apart. If curved surface of the drop is at right angles to the plates as shown in the figure and force exerted by the drop on a plate is $\dfrac{\pi\sigma\mathrm D}{\mathrm n}$, then find `n'\par
\qfig{p2_phy_s4_q06_1}
\qn{7} A parallel beam of light is incident from air at an angle `$\alpha$' on the side PQ of a right angled triangular prism of refractive index n = $\sqrt2$. Light undergoes total internal reflection in the prism at the face PR when $\alpha$ has a minimum value of 45$^\circ$. The angle `$\theta$' of the prism is `n' times 3$^\circ$. Find `n'.\par
\qfig{p2_phy_s4_q07_1}
\qn{8} A wire having a length L and cross-sectional area A is suspended at one of its ends from a ceiling. Density and Young's modulus of material of the wire are $\rho$ and Y, respectively. Its strain energy due to its own weight is $\dfrac{\rho^2\mathrm g^2\mathrm{AL}^3}{\mathrm{pY}}$. Calculate the value of p.
\qn{9} In a hypothetical system a particle of mass m and charge $-$3q is moving around a very heavy particle having charge q. Assuming Bohr's model to be applicable to this system, the orbital velocity of mass m when it is nearest to heavy particle is $\dfrac{3\mathrm q^2}{\mathrm x\varepsilon_0\mathrm h}$. Find the value of x.
\qn{10} A cylinder of radius r = 0.1 m and mass M = 2 kg is placed such that it is in contact simultaneously with a vertical and a horizontal surface as shown in figure. The coefficient of static friction is $\mu$ = $\left(\dfrac13\right)$ for both the surfaces. Find the distance d, in cm, from the centre of the cylinder at which a force F = 40 N should be applied so that the cylinder just starts rotating in the anticlockwise direction. Take g = 10 ms$^{-2}$.\par
\qfigi{p2_phy_s4_q10_1}
\raggedcolumns
\papersec{2}
\partsec{2}{CHEMISTRY}
\sect{SECTION-I(i)}
\qn{1} Consider the following unit cell of a metal:\par
\qfig{p2_chem_s1_q01_1}
Also consider the following shaded planes P and Q representing two 2D lattices.\par
\noindent P \qfigi{p2_chem_s1_q01_2}\quad Q \qfigi{p2_chem_s1_q01_3}\par
The ratio of packing fraction of P to Q is
\begin{qoptions}
\opt{A}{1.00}
\opt{B}{1.63}
\opt{C}{0.55}
\opt{D}{0.74}
\end{qoptions}
\qfullbegin
\qn{2} Vanillin (L), C$_8$H$_8$O$_3$, is isolated from vanilla beans. It give an intense colour with FeCl$_3$ and positive Tollen's test. It is not steam-distilled and does not react with HCl. It goes through following steps :\par
\qfig{p2_chem_s1_q02_0}
\qfullsplit
Which of the following is correct statement about various products?
\begin{qoptions}
\opt{A}{L is \qfigi{p2_chem_s1_q02_AL} and P is \qfigi{p2_chem_s1_q02_AP}}
\opt{B}{L is \qfigi{p2_chem_s1_q02_BL} and P is \qfigi{p2_chem_s1_q02_BP}}
\opt{C}{L is \qfigi{p2_chem_s1_q02_CL} and P is \qfigi{p2_chem_s1_q02_CP}}
\opt{D}{L is \qfigi{p2_chem_s1_q02_DL} and P is \qfigi{p2_chem_s1_q02_DP}}
\end{qoptions}
\qfullend
\qn{3} Predict the major product (P) in the following reaction :\par
\qfigi{p2_chem_s1_q03_1}
\begin{qoptions}
\opt{A}{\qfigi{p2_chem_s1_q03_A}}
\opt{B}{\qfigi{p2_chem_s1_q03_B}}
\opt{C}{\qfigi{p2_chem_s1_q03_C}}
\opt{D}{\qfigi{p2_chem_s1_q03_D}}
\end{qoptions}
\qn{4} The correct molecular orbital diagram for O$_2$ molecule in the ground state is
\begin{qoptions}
\opt{A}{\qfigi{p2_chem_s1_q04_A}}
\opt{B}{\qfigi{p2_chem_s1_q04_B}}
\opt{C}{\qfigi{p2_chem_s1_q04_C}}
\opt{D}{\qfigi{p2_chem_s1_q04_D}}
\end{qoptions}
\sect{SECTION-I(ii)}
\qn{5} Consider one mole of a real gas A, with equation of state P(V $-$ b) = RT where `b' is a constant. Which of the following is/are correct statement(s) about A?\par
(P = pressure, V = volume, T = temperature in Kelvin, R = 8.314 J K mol$^{-1}$)
\begin{qoptions}
\opt{A}{Density of A is less than ideal gas at all pressure values.}
\opt{B}{Potential energy of A, due intermolecular forces, is either greater than (or) equal to that of ideal gas}
\opt{C}{If 5J of heat is given to A, and 2J work is done by the A then, its internal energy increases by 3J, which is true for ideal gas also}
\opt{D}{Compressibility factor (Z) of A is directly proportional to `P'.}
\end{qoptions}
\qn{6} When copper is reacted with dil. HNO$_3$ a gas is liberated. Which of the following are true for the gas?
\begin{qoptions}
\opt{A}{Has 11 valence e's, where extra electron occupies an antibonding $\pi$ *2p orbital}
\opt{B}{An acid anhydride}
\opt{C}{Paramagnetic}
\opt{D}{Neutral oxide}
\end{qoptions}
\qn{7} The complex(es), which can exhibit the type of isomerism shown by [Co(en)$_3$]Cl$_3$, is(are)\par
[en = H$_2$NCH$_2$CH$_2$NH$_2$]
\begin{qoptions}
\opt{A}{[Pt(en)(SCN)$_2$]}
\opt{B}{[Zn(gly)$_2$]}
\opt{C}{[Co(NH$_3$)$_4$Cl$_2$]$^+$}
\opt{D}{[Cr(en)$_2$(H$_2$O)(SO$_4$)]$^+$}
\end{qoptions}
\sect{SECTION-II(i)}
\paraq{1}{2}
Consider the following process of folding/unfolding of a tripeptide segment in a large protein. The free energy change ($\Delta$G) will depend on the interaction of the unfolded tripeptide with the solvent (water) and with the rest of the protein in the folded state (see below). Assume that the tripeptide is made up of one non-polar (hydrophobic, shaded) and two polar (hydrophilic; unshaded) residues. Use the following $\Delta$G data of formation of interaction between residue with water or protein :\par
i) a non-polar residue and the solvent (water) $\Delta$G = +8 kJ mol$^{-1}$\par
ii) a non-polar residue and the rest of the protein $\Delta$G = $-$4 kJ mol$^{-1}$\par
iii) a polar residue and the solvent (water) $\Delta$G = $-$16 kJ mol$^{-1}$\par
iv) a polar residue and the rest of the protein $\Delta$G = $-$14 kJ mol$^{-1}$\par
\qfig{p2_chem_s3_para1_2_1}
\qn{1} What is the magnitude of $\Delta$G for the folding of the tripeptide segment in KJ/mol?
\qn{2} What is the magnitude of $\Delta$G for the folding of the tripeptide segment if all the three residues are polar in KJ/mol?
\paraq{3}{4}
Given below is synthesis of nylon-66. Assume 100 \% yield in every step and none of the organic compounds involved are in excess.\par
\qfig{p2_chem_s3_para3_4_1}
\qn{3} The value of x is \underline{\hspace{1.5cm}}.
\qn{4} The value of y is \underline{\hspace{1.5cm}}.
\sect{SECTION-II(ii)}
\qn{5} A solution of 0.1 M CH$_3$COOH is placed between parallel electrodes of cross-section area 4 cm$^2$, separated by 2 cm. For this solution, resistance measured is 100 $\Omega$. The elevation in boiling point of the 0.1 M CH$_3$COOH solution is x $\times$ 10$^{-5}$. The value of x is (use following information)\par
K$_\mathrm b$ = 0.5K kg / mol; $\Lambda^\infty_\mathrm m$ (H$^+$) = 300 S cm$^2$ mole$^{-1}$; $\Lambda^\infty_\mathrm m$ (CH$_3$COO$^-$) = 100 S cm$^2$ mole$^{-1}$
\qn{6} H$_3$C $-$ C(=O) $-$ CH$_3$ $+$ I$_2$ $\xrightarrow{\text{NaOH}}$ Iodoform reaction\par
How many molecules of I$_2$ are required for 3 molecules of CH$_3$COCH$_3$?
\qn{7} Find total number of inner orbital diamagnetic complexes :\par
[Fe(CN)$_6$]$^{3-}$, [Co(H$_2$O)$_6$]$^{2+}$, [Cr(en)$_3$]$^{3+}$, [Fe(CN)$_6$]$^{4-}$, [Co(H$_2$O)$_6$]$^{3+}$, [Cr(CN)$_6$]$^{3-}$, [PtCl$_6$]$^{2-}$, [Ni(H$_2$O)$_6$]$^{2+}$, [Cr(CO)$_6$]
\qn{8} Consider following statements :\par
(1) CO$_2$(g) $+$ C(s) $\longrightarrow$ 2CO(g), it is endothermic reaction,\par
(2) In copper extraction from CuFeS$_2$, CaSiO$_3$ is main slag formed.\par
(3) Brass is an alloy consisting of Zn metal.\par
(4) Cu $+$ H$_2$SO$_4$(conc.) $\longrightarrow$ CuSO$_4$ $+$ H$_2$.\par
(5) In Hall-Heroult process, anode is corroded forming CO$_2$ gas only.\par
(6) Due to non-metallic impurites melting point of iron is raised.\par
(7) In extraction of copper from copper pyrite, roasting process is employed.\par
How many of the following are correct statements.
\qn{9} Consider the following precipitate/s\par
PbS , CuS , HgS , CdS , As$_2$S$_3$, As$_2$S$_5$, SnS, SnS$_2$, Sb$_2$S$_3$, Sb$_2$S$_5$\par
x = Total number of yellow ppt/s which soluble in yellow ammonium sulphide\par
y = Total number of black ppt/s which insoluble in hot \& dil. HNO$_3$\par
Find the value of $\dfrac{\mathrm x}{\mathrm y}$.
\qn{10} \qfigi{p2_chem_s4_q09_1}\par
Identify total number of stereoisomer's if double bond is present in cis configuration in given compound (X)
\raggedcolumns
\papersec{2}
\partsec{3}{MATHEMATICS}
\sect{SECTION-I(i)}
\qn{1} Consider f : [0, $\infty$) $\to$ [0, $\infty$) is bijective function and it satisfies f(x)$\cdot$e$^{\mathrm{f(x)}}$ = x, then the value of\par
$\displaystyle\lim_{x\to\infty}\dfrac{\mathrm f(x)}{\ln x}$ is
\begin{qoptions}
\opt{A}{0}
\opt{B}{1}
\opt{C}{e}
\opt{D}{Does not exist}
\end{qoptions}
\qn{2} Real number x, y, z satisfies\par
x + xy + xyz = 1\par
y + yz + xyz = 2\par
z + xz + xyz = 4\par
If the largest possible value of xyz is $\dfrac{\mathrm a+\sqrt{\mathrm b}}{\mathrm c}$, where b \& c are prime numbers, then (a + b + c) equals
\begin{qoptions}
\opt{A}{20}
\opt{B}{22}
\opt{C}{24}
\opt{D}{26}
\end{qoptions}
\qn{3} If y is implicit differentiable function of x such that y(x + y)$^2$ = x, then $\displaystyle\int\dfrac{dx}{x+3y}$ is equal to
\begin{qoptions}
\opt{A}{$\dfrac12\ln\left((x+y)^2+1\right)+\mathrm C$}
\opt{B}{$\ln\left((x+y)^2+1\right)+\mathrm C$}
\opt{C}{$\dfrac12\ln\left((x-y)^2+1\right)+\mathrm C$}
\opt{D}{$\ln\left((x-y)^2+1\right)+\mathrm C$}
\end{qoptions}
Where C is constant of integration
\qn{4} Let A = \{1,2, 3, 4, 5, 6, 7, 8\} be a set. Let f : A $\to$ A be a function such that f (x) $\neq$ x $\forall$ x $\in$ A, f (x) is even whenever x is even and f(x) is odd whenever x is odd. Then total number of onto functions satisfying these conditions will be
\begin{qoptions}
\opt{A}{36}
\opt{B}{49}
\opt{C}{64}
\opt{D}{81}
\end{qoptions}
\sect{SECTION-I(ii)}
\qn{5} Let f : R $\to$ R, g : R $\to$ R and h : R $\to$ R be differentiable function such that f(x) = x$^3$ + 5x + 1, g(f(x)) = x and h(g(g(x))) = x $\forall$ x $\in$ R then
\begin{qoptions}
\opt{A}{h(0) = 7}
\opt{B}{h$'$(1) = 1216}
\opt{C}{g$'$(1) = $\dfrac15$}
\opt{D}{h(g(2)) = 19}
\end{qoptions}
\qn{6} There are 20 points lying on a circle, distance between two consecutive points is same. Three points are selected at random. X represent no. of ways to select them. Y represent no. of ways to select them such that no two adjacent objects are selected and Z represent no. of ways to select such that chosen objects are not diametrically opposite then
\begin{qoptions}
\opt{A}{X = 1140}
\opt{B}{Y = 800}
\opt{C}{Z = 960}
\opt{D}{Y $\cap$ Z = 660}
\end{qoptions}
\qn{7} An ellipse having eccentricity $\dfrac34$ passes through the point P (5, 4). If the nearer focus is (1,1) and equation of tangent at P on the ellipse 4x + 3y $-$ 32 = 0, then
\begin{qoptions}
\opt{A}{Length of latus rectum of ellipse is 17.5}
\opt{B}{Length of chord through (1,1) parallel to the tangent at P is 17.5}
\opt{C}{Length of perpendicular from the focus other than (1, 1) on the tangent at P is 35}
\opt{D}{Locus of point of intersection of perpendicular tangents drawn to ellipse is (x + 1)$^2$ + (y + 8)$^2$ = 575}
\end{qoptions}
\sect{SECTION-II(i)}
\paraq{1}{2}
Let f(x) is strictly increasing and continuous function in [0, 1] and for x $\in$ [0, 1] g(x) = $\displaystyle\int_0^x$ f (t) dt, also\par
$\displaystyle\int_0^1$ (f (t) + $|$f (t) $|$) dt = 6. Local minimum value of g(x) is $-$1 occurs at x = a, where a $\in$ (0, 1), then
\qn{1} $\displaystyle\int_0^1$ g (x) f (x) dx is equal to
\qn{2} $\displaystyle\int_a^1$ f (x) dx is equal to
\paraq{3}{4}
A tetrahedron is formed by the planes P$_1$ : y + z = 0. P$_2$ : z + x = 0, P$_3$ : x + y = 0 and P$_4$ : x + y + z = 3.
\qn{3} The volume V of the tetrahedron is :
\qn{4} Square of distance of P$_4$ from the vertex opposite to it is
\sect{SECTION-II(ii)}
\qn{5} Consider f(x) = $\dfrac{x^2+4}{|x|}$ and g(x) = p$|$x$|$ + 2, p $>$ 0. If the exhaustive set of values of `p' for which f(x) $>$ g(x) $\forall$ x $\in$ R $-$ \{0\} is (a, b) then a + 4b is
\qn{6} A fair coin is tossed 10 times and the outcomes are listed. Let E$_i$ be the event that the i$^{\mathrm{th}}$ outcome is a head and E$_m$ be the event that the list contains exactly m heads. If E$_i$ and E$_m$ are independent, then `m' is equal to
\qn{7} A perpendicular is drawn from a point on the line $\dfrac{x-1}2=\dfrac{y+1}{-1}=\dfrac z1$ to the plane x + y + z = 3 such that the foot of the perpendicular Q also lies on the plane x $-$ y + z = 3, if the coordinate of Q is (a, b, c) then $|$a $-$ b $-$ c$|$ is
\qn{8} If $\vec{\mathrm a}_1$ = a$\hat{\mathrm i}$ + b$\hat{\mathrm j}$, $\vec{\mathrm a}_2$ = $\vec{\mathrm a}_1$ $-$ $\left(2\hat{\mathrm i}+\sqrt5\hat{\mathrm j}\right)$ and $\vec{\mathrm a}_3$ = $\vec{\mathrm a}_1$ + $\left(2\hat{\mathrm i}-\sqrt5\hat{\mathrm j}\right)$, then\par
$\begin{vmatrix}1+\vec a_1\cdot\vec a_1&1+\vec a_1\cdot\vec a_2&1+\vec a_1\cdot\vec a_3\\1+\vec a_2\cdot\vec a_1&1+\vec a_2\cdot\vec a_2&1+\vec a_2\cdot\vec a_3\\1+\vec a_3\cdot\vec a_1&1+\vec a_3\cdot\vec a_2&1+\vec a_3\cdot\vec a_3\end{vmatrix}$ is equal to
\qn{9} If $\displaystyle\sum_{r=0}^{3n}a_r(x-4)^r=\sum_{r=0}^{3n}A_r(x-5)^r$ and a$_k$ = 1 $\forall$ k $\geqslant$ 2n and $\displaystyle\sum_{r=0}^{3n}d_r(x-8)^r=\sum_{r=0}^{3n}B_r(x-9)^r$ and\par
$\displaystyle\sum_{r=0}^{3n}d_r(x-12)^r=\sum_{r=0}^{3n}D_r(x-13)^r$ and d$_k$ = 1 $\forall$ k $\geqslant$ 2n, then the value of $\dfrac{\mathrm A_{2n}+\mathrm D_{2n}}{\mathrm B_{2n}}$.
\qn{10} If $\theta$ is the angle between $\bar a$ and $\bar b$ such that $\sqrt3(\bar a\cdot\bar b)+\left|\bar b-(\bar a\cdot\bar b)\bar a\right|=1+\bar b\cdot\bar b$ where $|\bar a|$ = 1 then cosec$^2\theta$ is
\end{multicols}
\end{document}
